% ======================================================================
% Finite-time dynamics: tabular strata and two-channel model
% Self-contained fragment (no preamble). Assumes amsmath, amssymb, amsthm,
% the macro \E (expectation symbol), and the environments
% theorem, proposition, lemma, corollary, remark, definition, proof.
% KL divergences are written \mathrm{KL}(\cdot\,\|\,\cdot) explicitly.
% v1-correctif-1: corrections after the counter-reading by a fresh Opus
%   (papiers/D/audits/contre-lecture-dynamique-v1/), written by the second daughter
%   session (Opus), NOT by Fable. STATUS: A REVOIR PAR FABLE. Not for diffusion.
%   New citation keys (to verify on PDF, R7): kakade2001natural, hopfield1984,
%   wilson1972, cohen1983, hofbauer2002, leonardos2022, wainwright2008, anthony2017,
%   zelikman2022, dong2023raft, gulcehre2023rest.
% ======================================================================

\section{Finite-time dynamics}\label{sec:dynamics}

The static results above describe Gibbs equilibria. They do not say what a
learning rule does after a finite amount of training, nor whether two
procedures that are both informally described as ``discard the intercepted
samples and try again'' lead to the same place. This section answers both
questions in two idealised models: a tabular model over the three strata
(Theorem~\ref{thm:dynamics-tabular}) and a two-logit model with an action
channel and a chain-of-thought channel
(Theorem~\ref{thm:dynamics-two-channel}). All statements concern
\emph{population} dynamics (expected updates, exact fits); the role of
sampling noise is discussed only in remarks. Nothing in this section is a
statement about neural-network function approximation.

\subsection{Setting and conventions}\label{sec:dynamics-setting}

Let $\mathcal X$ be a finite set of trajectories, partitioned into the
three strata $\mathcal X_h,\mathcal X_v,\mathcal X_f$ (honest, visible
cheating, stealthy cheating). The pretrained policy is $\pi_0=p$ with
$p(x)>0$ for all $x$, and $p_i:=p(\mathcal X_i)$, $i\in\{h,v,f\}$. As in the
static setting, the task reward $R$ and the survival probability $\beta$
are constant on strata: $R=0$ on $\mathcal X_h$, $R=\Delta R>0$ on
$\mathcal X_v\cup\mathcal X_f$, and
$0\le\beta_v<\beta_f\le\beta_h\le 1$, $\beta_f>0$; in particular every stratum has positive mass, $p_i>0$, and penalties satisfy $C\ge0$. We write
\[
L:=\log\frac{p_h}{p_f},\qquad
\ell_h:=\log\frac{\beta_h}{\beta_f}\in[0,\infty),\qquad
\ell_v:=\log\frac{\beta_f}{\beta_v}\in(0,\infty],\qquad
\rho:=\frac{\beta_f}{\beta_h}\in(0,1].
\]
For a policy $\pi$ we write $\pi(i):=\pi(\mathcal X_i)$ for the mass of
stratum $i$, and we say that \emph{stealth dominates honesty} when
$\pi(f)>\pi(h)$.

\paragraph{Tabular softmax and natural gradient.}
The policy is $\pi_\theta(x)=e^{\theta_x}/\sum_{y}e^{\theta_y}$,
$\theta\in\mathbb R^{\mathcal X}$, initialised at $\theta(0)=\log p$, so
that $\pi_{\theta(0)}=\pi_0$. Its Fisher information matrix is
\[
F(\theta)=\E_{x\sim\pi_\theta}\bigl[\nabla_\theta\log\pi_\theta(x)\,
\nabla_\theta\log\pi_\theta(x)^{\top}\bigr]
=\operatorname{diag}(\pi_\theta)-\pi_\theta\pi_\theta^{\top}.
\]
For a
function $V$ that is $C^1$ on a neighbourhood of the open simplex, the
\emph{natural gradient flow} of $V$ is
$\dot\theta=F(\theta)^{\dagger}\nabla_\theta (V\circ\pi_\theta)$ and the
\emph{natural gradient ascent} with step $\gamma>0$ is
$\theta_{k+1}=\theta_k+\gamma F(\theta_k)^{\dagger}\nabla_\theta
(V\circ\pi_{\theta})(\theta_k)$, where $F^\dagger$ is the Moore--Penrose
pseudo-inverse. We denote by $\mathbf 1$ the all-ones vector.

\begin{lemma}[Natural gradient in the tabular softmax family]
\label{lem:npg-softmax}
Let $V$ be $C^1$ on a neighbourhood of the open simplex and let
$G(\pi):=\nabla_\pi V(\pi)\in\mathbb R^{\mathcal X}$. Then
\[
F(\theta)^{\dagger}\nabla_\theta (V\circ\pi_\theta)
= G(\pi_\theta)-\Bigl(\tfrac{1}{|\mathcal X|}\textstyle\sum_y G_y(\pi_\theta)\Bigr)\mathbf 1 .
\]
Consequently, along the natural gradient flow,
$\frac{d}{dt}\log\pi_t(x)=G_x(\pi_t)-\E_{\pi_t}[G(\pi_t)]$ and
$\frac{d}{dt}V(\pi_t)=\sum_x\pi_t(x)\bigl(G_x(\pi_t)-\E_{\pi_t}[G(\pi_t)]\bigr)^2\ge 0$;
along natural gradient ascent,
$\pi_{k+1}(x)\propto\pi_k(x)\exp\bigl(\gamma\,G_x(\pi_k)\bigr)$.
\end{lemma}

\begin{proof}
Since $\partial\pi_\theta(y)/\partial\theta_x=\pi_\theta(y)(\mathbf 1\{x=y\}-\pi_\theta(x))$,
the chain rule gives
$\nabla_\theta(V\circ\pi_\theta)=(\operatorname{diag}(\pi_\theta)-\pi_\theta\pi_\theta^\top)G=F G$.
The matrix $F$ is symmetric, and $(Fz)_x=\pi_\theta(x)(z_x-\E_{\pi_\theta}[z])$
vanishes for all $x$ if and only if $z$ is constant, because $\pi_\theta(x)>0$;
hence $\ker F=\operatorname{span}(\mathbf 1)$ and $F^\dagger F$ is the
orthogonal projector $\Pi$ onto $\mathbf 1^{\perp}$, which gives the first
identity. The policy $\pi_\theta$ is invariant under
$\theta\mapsto\theta+c\mathbf 1$, and
$\frac{d}{dt}\log\pi_\theta(x)=\dot\theta_x-\sum_y\pi_\theta(y)\dot\theta_y$,
which gives the second identity; the third follows from
$\frac{d}{dt}V=(FG)^\top\Pi G=G^\top F G$ (as $F\Pi=F$) and
$G^\top FG=\sum_x\pi(x)(G_x-\E_\pi G)^2$. The discrete statement is the
definition of the softmax applied to $\theta_{k+1}=\theta_k+\gamma\Pi G(\pi_k)$.
\end{proof}

\begin{definition}[Gibbs family]\label{def:gibbs-family}
For $r:\mathcal X\to\mathbb R$ and $\zeta\ge 0$, let
$\mathcal G_\zeta[r](x):=p(x)e^{\zeta r(x)}/\sum_y p(y)e^{\zeta r(y)}$.
If $r$ is constant on strata, with values $r_h,r_v,r_f$, then
$\mathcal G_\zeta[r](i)\propto p_i e^{\zeta r_i}$.
\end{definition}

The three static operators are members of this family at $\zeta=\eta$:
$\pi_{\mathrm{no}}=\mathcal G_\eta[R]$,
$\pi_B=\mathcal G_\eta[\bar r_C]$ with
$\bar r_C=\beta R-(1-\beta)C$, and, when $\beta_v>0$,
$\pi_A=\mathcal G_\eta[R+\eta^{-1}\log\beta]$.

\paragraph{Convention for $\beta_v=0$ under operator~A.}
When $\beta_v=0$ the effective reward $R+\eta^{-1}\log\beta$ equals
$-\infty$ on $\mathcal X_v$. We then define every object attached to
operator~A on the reduced space $\mathcal X_h\cup\mathcal X_f$ with $p$
renormalised, and set the mass of $\mathcal X_v$ to zero for $t>0$. This is
the pointwise limit, for each fixed $t>0$, of the formulas below as
$\beta_v\downarrow 0$, and no ratio between the masses of $h$ and $f$ in
those formulas depends on $\beta_v$. With this convention $\beta_v^{\lambda}:=0$
for $\lambda>0$.

\subsection{Tabular dynamics: three learning rules}

We consider three learning rules. Only rules~(b) and~(c) discard intercepted samples; rule~(a) applied to operator~A is a penalty in $\log\beta$, not a rejection procedure.

\smallskip\noindent\textbf{(a) Regularised natural policy gradient.}
For an effective reward $r:\mathcal X\to\mathbb R$, the rule is natural
gradient flow or ascent on
$V_r(\pi):=\E_\pi[r]-\eta^{-1}\mathrm{KL}(\pi\,\|\,\pi_0)$. The cases of
interest are $r=R$ (no monitor), $r=\bar r_C$ (operator~B) and
$r=R+\eta^{-1}\log\beta$ (operator~A, in its variational form). For operator~A this rule discards nothing and resamples nothing: it ascends a reward augmented by the bonus $\eta^{-1}\log\beta(x)$, which requires the survival probability of each trajectory (a binary interception gives no unbiased estimate of $\log\beta(x)$) and is an infinite penalty when $\beta_v=0$. It is the dynamic counterpart of the modelling choice made for the static operator~A, and the statements derived from it (Theorem~\ref{thm:dynamics-tabular}(a2), Corollary~\ref{cor:hitting-times}(iii), Proposition~\ref{prop:equal-time}) are properties of this relaxation.

\smallskip\noindent\textbf{(b) Iterated rejection fine-tuning.}
At round $k$, trajectories are sampled from $\pi_k$, intercepted ones are
discarded, the accepted ones are reweighted (or subsampled) with a
selection weight $\omega(R)$, where $\omega:\{0,\Delta R\}\to(0,\infty)$ is
nondecreasing, and $\pi_{k+1}$ is fitted exactly to the resulting
distribution:
\begin{equation}\label{eq:rejection-ft}
\pi_{k+1}(x)\;\propto\;\pi_k(x)\,\beta(x)\,\omega(R(x)),\qquad \pi_0=p .
\end{equation}
We also consider the anchored variant, for $\alpha\in(0,1)$,
\begin{equation}\label{eq:rejection-ft-anchored}
\pi_{k+1}:=\arg\min_{\pi}\Bigl\{(1-\alpha)\,\mathrm{KL}(\pi\,\|\,\kappa_k)
+\alpha\,\mathrm{KL}(\pi\,\|\,\pi_0)\Bigr\},
\qquad \kappa_k\propto\pi_k\,\beta\,\omega(R),
\end{equation}
the minimum being over distributions on the support of $\kappa_k$. With $\omega(R)=e^{\gamma R}$ the unanchored rule~\eqref{eq:rejection-ft} is the repeated rejection already discussed for the static model (remark on operator~A as a modelling choice, $\pi_T\propto\beta^T e^{T\eta_1R}\pi_0$); part~(b) below generalises it to arbitrary $\omega$, the anchored variant and best-of-$N$ selection. Rule~(b) with $\omega=e^{\gamma R}$ is also the filtering iteration of expert iteration and self-training on filtered samples~\cite{anthony2017,zelikman2022,dong2023raft,gulcehre2023rest}.

\smallskip\noindent\textbf{(c) Discard policy gradient.}
Trajectories are drawn from $\pi_\theta$ until $N\ge 2$ have been accepted;
call them $x_1,\dots,x_N$. The update direction is the REINFORCE estimator
computed on accepted samples only, with the accepted-sample mean as
baseline and the score of the \emph{sampling} policy:
\begin{equation}\label{eq:discard-pg-estimator}
\hat g:=\frac1N\sum_{n=1}^N\bigl(R(x_n)-\bar R_N\bigr)\nabla_\theta\log\pi_\theta(x_n),
\qquad \bar R_N:=\frac1N\sum_{n=1}^N R(x_n).
\end{equation}
For a policy $\pi$ with $\E_\pi[\beta]>0$ define
\begin{equation}\label{eq:accepted-mean}
\bar Z(\pi):=\E_\pi[\beta],\qquad
b(\pi):=\frac{\E_\pi[\beta R]}{\E_\pi[\beta]},\qquad
a_x(\pi):=\frac{\beta(x)\bigl(R(x)-b(\pi)\bigr)}{\bar Z(\pi)} .
\end{equation}
Thus $b(\pi)$ is the mean reward conditional on acceptance, and
$\E_\pi[a(\pi)]=0$.

\paragraph{Geometry of rule (c).} Throughout, rule~(c) means the natural-gradient version: the parameter moves along $F(\theta)^{\dagger}\hat g$, where $F(\theta)$ is the Fisher information matrix of the sampling policy. By (c1) below its expected direction is $a(\pi_\theta)$ modulo $\mathbf 1$. Applied without this preconditioning, the estimator follows the plain gradient flow $\dot\theta=F(\theta)a(\pi_\theta)$, which does not close on strata when a stratum has several trajectories, and for which (c2)(ii)--(iii) are false: with $p=(0.6,0.3,0.1)$, $\beta=(1,0.2,0.5)$, $\Delta R=1$ and one trajectory per stratum, the natural flow converges to $(0,0.548,0.452)$, whereas the plain flow is at $(1.0\cdot10^{-6},\,0.9965,\,0.0035)$ at $t=10^{5}$ and still drifting, the quantity conserved by the natural flow moving from $0.754$ to $5.03$, and the stealthy share among cheaters falling below its unmonitored natural-flow value $0.25$.

\begin{theorem}[Tabular finite-time dynamics]\label{thm:dynamics-tabular}
In the setting of Section~\ref{sec:dynamics-setting}:

\smallskip\noindent\textup{(a) Time-slice identity.}
Let $r:\mathcal X\to\mathbb R$ be arbitrary. The natural gradient flow of
$V_r$ started at $\pi_0$ exists for all $t\ge 0$ and satisfies
\begin{equation}\label{eq:time-slice}
\pi_t=\mathcal G_{\eta\lambda(t)}[r],\qquad \lambda(t)=1-e^{-t/\eta}.
\end{equation}
Natural gradient ascent on $V_r$ with step $\gamma>0$ satisfies
$\pi_k=\mathcal G_{\eta\lambda_k}[r]$ with
$\lambda_k=1-(1-\gamma/\eta)^k$; the sequence $\lambda_k$ is nondecreasing
with values in $[0,1]$ if and only if $\gamma\le\eta$, and $\lambda_1=1$ when
$\gamma=\eta$. In particular:
\begin{enumerate}
\item[\textup{(a1)}] for $r=R$ and for $r=\bar r_C$, the trajectory is
exactly the corresponding \emph{static} family evaluated at the effective
pressure $\eta_{\mathrm{eff}}(t)=\eta\lambda(t)$;
\item[\textup{(a2)}] for operator~A,
$\pi_t(x)\propto p(x)\,\beta(x)^{\lambda(t)}\,e^{\eta\lambda(t)R(x)}$: the
pressure is $\eta\lambda(t)$ but the filter enters with exponent
$\lambda(t)<1$, so $\pi_t$ is \emph{not} the static operator~A at pressure
$\eta_{\mathrm{eff}}(t)$.
\end{enumerate}

\smallskip\noindent\textup{(b) Iterated rejection fine-tuning.}
Under \eqref{eq:rejection-ft}, $\pi_k\propto p\,\beta^{k}\,\omega(R)^k$, so
that for every $k\ge 0$
\begin{equation}\label{eq:rejection-logodds}
\log\frac{\pi_k(f)}{\pi_k(h)}=-L+k\,\vartheta,\qquad
\vartheta:=\log\frac{\omega(\Delta R)}{\omega(0)}-\ell_h,
\qquad
\log\frac{\pi_k(f)}{\pi_k(v)}=\log\frac{p_f}{p_v}+k\,\ell_v
\end{equation}
(the last quantity being $+\infty$ for $k\ge1$ if $\beta_v=0$). Under the
anchored rule \eqref{eq:rejection-ft-anchored}, the same formulas hold with
$k$ replaced by $c_k:=\frac{1-\alpha}{\alpha}\bigl(1-(1-\alpha)^k\bigr)$,
which increases to $c_\infty=(1-\alpha)/\alpha$.

\smallskip\noindent\textup{(c) Discard policy gradient.}
\begin{enumerate}
\item[\textup{(c1)}] \emph{Identification.} For every $\theta$,
$\E[\hat g]=\frac{N-1}{N}\,\nabla_\theta\bigl(b\circ\pi_\theta\bigr)
=\frac{N-1}{N}\,F(\theta)\,a(\pi_\theta)$ and $\nabla_\pi b(\pi)=a(\pi)$.
Hence the expected discard policy gradient is, up to the factor $(N-1)/N$,
the gradient of the accepted-sample mean reward $b$, and its natural
gradient direction is $a(\pi_\theta)$ modulo $\mathbf 1$.
\item[\textup{(c2)}] \emph{Unregularised flow.} Let $\pi_t$ be the natural
gradient flow of $b$ started at $\pi_0$ (that is, $\dot\theta=a(\pi_\theta)$
modulo $\mathbf 1$). It exists for all $t\ge0$, and, writing
$m:=\beta_v\pi_t(v)+\beta_f\pi_t(f)$ and $\bar Z=\beta_h\pi_t(h)+m$,
\begin{equation}\label{eq:discard-a}
\frac{d}{dt}\log\pi_t(h)=-\frac{\beta_h\,m\,\Delta R}{\bar Z^{2}},\quad
\frac{d}{dt}\log\pi_t(v)=\frac{\beta_v\beta_h\,\pi_t(h)\,\Delta R}{\bar Z^{2}},\quad
\frac{d}{dt}\log\pi_t(f)=\frac{\beta_f\beta_h\,\pi_t(h)\,\Delta R}{\bar Z^{2}} ,
\end{equation}
the right-hand sides being $a_h(\pi_t)$, $a_v(\pi_t)$ and $a_f(\pi_t)$.
Consequently:
\begin{enumerate}
\item[\textup{(i)}] $\displaystyle\frac{d}{dt}\log\frac{\pi_t(f)}{\pi_t(h)}
=\frac{\beta_f\Delta R+(\beta_h-\beta_f)\,b(\pi_t)}{\bar Z(\pi_t)}
\;\ge\;\rho\,\Delta R>0$. Hence
$\pi_t(h)\le (p_h/p_f)\,e^{-\rho\Delta R\,t}$, and if $p_h>p_f$ stealth
dominates honesty from a time $\tau\le L/(\rho\,\Delta R)$ onwards, for
\emph{every} $\beta_f>0$.
\item[\textup{(ii)}] $\displaystyle\frac{d}{dt}\log\frac{\pi_t(f)}{\pi_t(v)}
=\frac{(\beta_f-\beta_v)\beta_h\,\pi_t(h)\,\Delta R}{\bar Z(\pi_t)^{2}}>0$,
and the quantity $\pi_t(v)\,\pi_t(f)^{-\beta_v/\beta_f}$ is conserved:
$\pi_t(v)=p_v\bigl(\pi_t(f)/p_f\bigr)^{\beta_v/\beta_f}$ for all $t$.
\item[\textup{(iii)}] $\pi_t$ converges to $(\pi_\infty(h),\pi_\infty(v),\pi_\infty(f))=(0,1-x^\ast,x^\ast)$,
where $x^\ast$ is the unique solution in $(p_f,1)$ of
$x+p_v(x/p_f)^{\beta_v/\beta_f}=1$. Moreover
$x^\ast>p_f/(p_f+p_v)$, the limit of the same flow without monitor
($\beta\equiv1$); and $x^\ast=1-p_v$ when $\beta_v=0$.
\end{enumerate}
\end{enumerate}
\end{theorem}

\begin{proof}
\emph{(a)} We have
$\nabla_\pi V_r(\pi)=r-\eta^{-1}\bigl(\log(\pi/p)+\mathbf 1\bigr)$, and
$\log(\pi_\theta/p)=\theta-\log p$ modulo $\mathbf 1$. By
Lemma~\ref{lem:npg-softmax} the flow is
$\dot\theta=\Pi\bigl(r-\eta^{-1}(\theta-\log p)\bigr)$, a linear equation,
hence globally solvable. Put $\phi:=\theta-\log p$; then
$\dot\phi=r-\phi/\eta+c(t)\mathbf 1$ for a continuous scalar function $c$,
with $\phi(0)=0$, and variation of constants gives
$\phi(t)=\eta(1-e^{-t/\eta})\,r+\bigl(\int_0^te^{-(t-s)/\eta}c(s)\,ds\bigr)\mathbf 1$.
Since $\pi_t\propto p\,e^{\phi(t)}$, this is \eqref{eq:time-slice}. In
discrete time, $\phi_{k+1}=(1-\gamma/\eta)\phi_k+\gamma r$ modulo
$\mathbf 1$, and $\phi_k=\eta\lambda_kr$ modulo $\mathbf 1$ follows by
induction from $\eta\lambda_{k+1}=(1-\gamma/\eta)\eta\lambda_k+\gamma$. The
monotonicity statement is immediate. Items (a1) and (a2) are
\eqref{eq:time-slice} for the stated $r$; for (a2),
$\exp(\eta\lambda\cdot\eta^{-1}\log\beta)=\beta^{\lambda}$, and the case
$\beta_v=0$ is covered by the convention above.

\emph{(b)} The first claim is an induction on \eqref{eq:rejection-ft}, and
\eqref{eq:rejection-logodds} follows by summing over strata, on which
$\beta$ and $R$ are constant (for $f$ against $v$ the factor $\omega$
cancels because $R$ takes the same value). For the anchored rule, the
objective in \eqref{eq:rejection-ft-anchored} equals
$\mathrm{KL}(\pi\,\|\,\mu_k)$ plus a constant, where
$\mu_k\propto\kappa_k^{1-\alpha}\pi_0^{\alpha}$, so
$\pi_{k+1}=\mu_k$. Writing $\psi:=\log\beta+\log\omega(R)$ and
$\phi_k:=\log(\pi_k/p)$ modulo constants, this reads
$\phi_{k+1}=(1-\alpha)(\phi_k+\psi)$, whence
$\phi_k=\psi\sum_{j=1}^k(1-\alpha)^j=c_k\psi$.

\emph{(c1)} Conditionally on acceptance, a draw from $\pi_\theta$ has law
$\pi^{\mathrm{acc}}(x)=\pi_\theta(x)\beta(x)/\bar Z$, and the $N$ accepted
samples are independent with this law. Let
$s_n:=\nabla_\theta\log\pi_\theta(x_n)=e_{x_n}-\pi_\theta$ and
$R_n:=R(x_n)$. By independence,
\[
\E\bigl[(R_n-\bar R_N)s_n\bigr]
=\Bigl(1-\tfrac1N\Bigr)\E[R_ns_n]-\tfrac1N\sum_{m\ne n}\E[R_m]\,\E[s_n]
=\tfrac{N-1}{N}\Bigl(\E[R_1s_1]-\E[R_1]\E[s_1]\Bigr),
\]
and the $y$-th coordinate of $\E[R_1s_1]-\E[R_1]\E[s_1]$ is the covariance,
under $\pi^{\mathrm{acc}}$, of $R$ and $\mathbf 1\{x=y\}$, namely
$\pi^{\mathrm{acc}}(y)(R(y)-b)=\pi_\theta(y)a_y(\pi_\theta)$. Thus
$\E[\hat g]=\frac{N-1}{N}\,\pi_\theta\circ a$ (coordinatewise product), and
$\pi_\theta\circ a=Fa$ because $\E_{\pi_\theta}[a]=0$. On the other hand the
quotient rule gives
$\partial b/\partial\pi(x)=\beta(x)R(x)/\bar Z-\E_\pi[\beta R]\beta(x)/\bar Z^{2}=a_x(\pi)$,
and Lemma~\ref{lem:npg-softmax} gives
$\nabla_\theta(b\circ\pi_\theta)=Fa$ and the natural gradient direction.

\emph{(c2)} The vector $a(\pi)$ is constant on strata and depends on $\pi$
only through the stratum masses, so by Lemma~\ref{lem:npg-softmax}
(with $\E_\pi[a]=0$) $\frac{d}{dt}\log\pi_t(x)=a_x(\pi_t)$ and the stratum
masses obey $\frac{d}{dt}\log\pi_t(i)=a_i(\pi_t)$. Inserting
$b=m\Delta R/\bar Z$ and $\Delta R-b=\beta_h\pi(h)\Delta R/\bar Z$ in
\eqref{eq:accepted-mean} gives \eqref{eq:discard-a}. In particular
$\pi_t(f)$ is nondecreasing, so $\bar Z\ge\beta_fp_f>0$ along the solution,
$a$ stays bounded, and the solution is global.

(i) Subtracting, $a_f-a_h=\bigl(\beta_f(\Delta R-b)+\beta_hb\bigr)/\bar Z$,
which is the stated expression. Since $b\ge0$, $\beta_h\ge\beta_f$ and
$\bar Z\le\max_x\beta(x)=\beta_h$, it is at least
$\beta_f\Delta R/\beta_h=\rho\Delta R$. Integrating,
$\pi_t(h)/\pi_t(f)\le(p_h/p_f)e^{-\rho\Delta Rt}$, which gives both
consequences.

(ii) The rate is $a_f-a_v$ read off \eqref{eq:discard-a}. Since
$a_v=(\beta_v/\beta_f)a_f$, the function
$\log\pi_t(v)-(\beta_v/\beta_f)\log\pi_t(f)$ has zero derivative.

(iii) By (i), $\pi_t(h)\to0$; $\pi_t(f)$ is nondecreasing and bounded, so it
converges to some $x^\ast$, and by (ii) $\pi_t(v)\to p_v(x^\ast/p_f)^{\beta_v/\beta_f}$.
The masses sum to one, which is the equation for $x^\ast$. Its left-hand
side is continuous and strictly increasing in $x$, equals $p_f+p_v<1$ at
$x=p_f$ and exceeds $1$ at $x=1$, so the solution is unique and lies in
$(p_f,1)$. At $x^{\mathrm{no}}:=p_f/(p_f+p_v)$ the left-hand side equals
$x^{\mathrm{no}}+p_v(p_f+p_v)^{-\beta_v/\beta_f}<x^{\mathrm{no}}+p_v(p_f+p_v)^{-1}=1$,
because $\beta_v/\beta_f<1$ and $p_f+p_v<1$; hence
$x^\ast>x^{\mathrm{no}}$. Without monitor, $a=R-b$ and the flow is
$\pi_t\propto p\,e^{tR}$, with limit mass $x^{\mathrm{no}}$ on $f$. The case $\beta_v=0$ is read off the equation.
\end{proof}

\begin{remark}[Equivalent forms of rule (a)]\label{rem:mirror}
The discrete rule in part~(a) is
$\pi_{k+1}\propto\pi_k^{\,1-\gamma/\eta}\,\pi_0^{\,\gamma/\eta}\,e^{\gamma r}$,
which is also the mirror-ascent step
$\pi_{k+1}=\arg\max_\pi\{\langle\nabla_\pi V_r(\pi_k),\pi\rangle-\gamma^{-1}\mathrm{KL}(\pi\,\|\,\pi_k)\}$.
This multiplicative form of regularised natural policy gradient, and its
linear convergence, are known \cite{cen2022fast}; part~(a) is the
observation that, started at the anchor $\pi_0$, the whole trajectory stays
inside the one-parameter family
$\{\mathcal G_\zeta[r]\}_{0\le\zeta\le\eta}$.
\end{remark}

\begin{corollary}[Exact dominance times]\label{cor:hitting-times}
Assume $p_h>p_f$, i.e.\ $L>0$. For $\Delta\in\mathbb R$ let
\[
\tau(\Delta):=\begin{cases}
-\eta\log\bigl(1-\dfrac{L}{\eta\Delta}\bigr) & \text{if }\eta\Delta>L,\\[1.5ex]
+\infty & \text{otherwise.}
\end{cases}
\]
\begin{enumerate}
\item[\textup{(i)}] Under rule \textup{(a)} with a reward $r$ constant on
strata and $\Delta:=r_f-r_h$, one has $\pi_t(f)>\pi_t(h)$ if and only if
$t>\tau(\Delta)$. In discrete time with $\gamma<\eta$, the first dominance
step is
$\bigl\lfloor\log(1-\tfrac{L}{\eta\Delta})/\log(1-\tfrac{\gamma}{\eta})\bigr\rfloor+1$
when $\eta\Delta>L$ (it is $1$ when $\gamma=\eta$), and dominance never
occurs otherwise.
\item[\textup{(ii)}] The relevant values are
$\Delta_{\mathrm{no}}=\Delta R$;
$\Delta_B=\beta_f\Delta R-(\beta_h-\beta_f)C$; and
$\Delta_A=\Delta R-\eta^{-1}\ell_h$ (also when $\beta_v=0$). Thus
\begin{align*}
\tau_A&=-\eta\log\Bigl(1-\frac{L}{\eta\Delta R-\ell_h}\Bigr)
&&\text{if }\ \eta\Delta R>L+\ell_h,\\
\tau_B&=-\eta\log\Bigl(1-\frac{L}{\eta\bigl(\beta_f\Delta R-(\beta_h-\beta_f)C\bigr)}\Bigr)
&&\text{if }\ \eta\bigl(\beta_f\Delta R-(\beta_h-\beta_f)C\bigr)>L,
\end{align*}
and $+\infty$ otherwise; the finiteness conditions are exactly the static
dominance thresholds.
\item[\textup{(iii)}] $\tau$ is nonincreasing in $\Delta$, so
$\tau_{\mathrm{no}}\le\tau_A$ and $\tau_{\mathrm{no}}\le\tau_B$. Moreover
$\tau(\Delta)\ge L/\Delta$ for $\Delta>0$, and
$\tau(\Delta)\to L/\Delta$ as $\eta\to\infty$ with $\Delta$ fixed. In
particular $\tau_A\to L/\Delta R=\lim\tau_{\mathrm{no}}$ as
$\eta\to\infty$: in this limit the filter of operator~A, read through the variational rule~(a) (a $\log\beta$ penalty, not a rejection), does not delay
dominance, whereas $\tau_B\to L/\Delta_B$ if $\Delta_B>0$ and
$\tau_B=+\infty$ for every $\eta$ if $\Delta_B\le0$.
\item[\textup{(iv)}] Under rule \textup{(b)}, stealth dominates honesty at
some round if and only if $\vartheta>0$, that is, if and only if
$\log(\omega(\Delta R)/\omega(0))>\ell_h$, and then exactly from
round $k^\ast=\lfloor L/\vartheta\rfloor+1$ onwards. Under the anchored rule
the condition is $\vartheta>0$ and $c_\infty\vartheta>L$, and the first
round is then
\[
\Bigl\lfloor\log\Bigl(1-\frac{\alpha L}{(1-\alpha)\vartheta}\Bigr)\Big/\log(1-\alpha)\Bigr\rfloor+1 .
\]
\item[\textup{(v)}] Under the unregularised rule \textup{(c)} (natural-gradient form), stealth
dominates honesty from a time $\tau\le L/(\rho\Delta R)$ onwards, whatever
$\beta_f>0$.
\end{enumerate}
\end{corollary}

\begin{proof}
(i) By Theorem~\ref{thm:dynamics-tabular}(a) and
Definition~\ref{def:gibbs-family},
$\log(\pi_t(f)/\pi_t(h))=-L+\eta\lambda(t)\Delta$. Since $L>0$ and
$\lambda$ increases from $0$ to $1$ (never reaching $1$), this is positive
for some $t$ if and only if $\Delta>0$ and $\eta\Delta>L$, and then exactly
for $\lambda(t)>L/(\eta\Delta)$, i.e.\ $t>\tau(\Delta)$. In discrete time
the condition $\lambda_k>L/(\eta\Delta)$ reads
$(1-\gamma/\eta)^k<1-L/(\eta\Delta)$, i.e.\
$k>\log(1-L/(\eta\Delta))/\log(1-\gamma/\eta)$. (ii) The three differences
are computed from the definitions:
$\bar r_C(f)-\bar r_C(h)=\beta_f\Delta R-(1-\beta_f)C+(1-\beta_h)C$, and
$(R+\eta^{-1}\log\beta)(f)-(R+\eta^{-1}\log\beta)(h)=\Delta R-\eta^{-1}\ell_h$.
(iii) Monotonicity is clear; $-\log(1-z)\ge z$ gives
$\tau(\Delta)\ge L/\Delta$, and $-\eta\log(1-L/(\eta\Delta))\to L/\Delta$ as
$\eta\to\infty$. For operator~A, $\eta\Delta_A=\eta\Delta R-\ell_h$, so
$\tau_A=-\eta\log(1-L/(\eta\Delta R-\ell_h))\to L/\Delta R$.
(iv) By \eqref{eq:rejection-logodds} the log-odds are $-L+k\vartheta$,
positive if and only if $\vartheta>0$ and $k>L/\vartheta$; in the anchored
case, $c_k>L/\vartheta$ reads
$(1-\alpha)^k<1-\alpha L/((1-\alpha)\vartheta)$. (v) is
Theorem~\ref{thm:dynamics-tabular}(c2)(i).
\end{proof}

\paragraph{Absolute amplification at equal training time.}
In the static setting, operator~A \emph{amplifies} stealth in absolute
terms when $\pi_A(f)>\pi_{\mathrm{no}}(f)$; a direct computation shows that
this holds if and only if
$(\beta_f-\beta_v)\,p_v\,e^{\eta\Delta R}>(\beta_h-\beta_f)\,p_h$, i.e.\
always when $\beta_h=\beta_f$, and otherwise if and only if
$\eta\Delta R>T_{\mathrm{amp}}:=\log\frac{p_h(\beta_h-\beta_f)}{p_v(\beta_f-\beta_v)}$.
The dynamic counterpart compares, at the \emph{same} time $t$ and the same
$\eta$, the run $\pi^{A}_t$ of rule~(a) for operator~A with the unmonitored
run $\pi^{\mathrm{no}}_t$ of rule~(a) for $r=R$. Both runs have the same
$\lambda(t)$.

\begin{proposition}[Equal-time amplification under operator~A]
\label{prop:equal-time}
Let $t>0$ and $\lambda=\lambda(t)\in(0,1)$.
\begin{enumerate}
\item[\textup{(i)}] \emph{Exact criterion.}
$\pi^{A}_t(f)>\pi^{\mathrm{no}}_t(f)$ if and only if
\begin{equation}\label{eq:equal-time-criterion}
p_v\bigl(\beta_f^{\lambda}-\beta_v^{\lambda}\bigr)\,e^{\eta\lambda\Delta R}
\;>\;p_h\bigl(\beta_h^{\lambda}-\beta_f^{\lambda}\bigr).
\end{equation}
\item[\textup{(ii)}] \emph{Explicit sufficient time.} If $\beta_h=\beta_f$,
\eqref{eq:equal-time-criterion} holds for all $t>0$. If $\beta_h>\beta_f$,
it holds at every $t>0$ such that
$\eta_{\mathrm{eff}}(t)\,\Delta R\ge T_{\mathrm{amp}}$; that is, for all
$t>0$ if $T_{\mathrm{amp}}\le0$, and for all
$t\ge-\eta\log\bigl(1-T_{\mathrm{amp}}/(\eta\Delta R)\bigr)$ if
$0<T_{\mathrm{amp}}<\eta\Delta R$.
\item[\textup{(iii)}] \emph{Short times.} Let $\beta_h>\beta_f$. If
$\beta_v=0$, or if $\beta_v>0$ and $p_v\ell_v>p_h\ell_h$, then
\eqref{eq:equal-time-criterion} holds for all sufficiently small $t>0$; if
$\beta_v>0$ and $p_v\ell_v<p_h\ell_h$, it fails for all sufficiently small
$t>0$.
\item[\textup{(iv)}] \emph{Structure.} Let $\beta_h>\beta_f$ and, for
$\lambda>0$, let $g(\lambda)$ be the logarithm of the ratio of the two sides
of \eqref{eq:equal-time-criterion}. Then $g$ is convex on $(0,\infty)$ if
$\beta_f^2\ge\beta_v\beta_h$ and concave if
$\beta_f^2\le\beta_v\beta_h$. Hence the set of $\lambda\in(0,1)$ at which
amplification holds is an interval in the concave case, and the complement
in $(0,1)$ of a relatively closed, possibly empty, subinterval in the convex case (the boundary case $p_v\ell_v=p_h\ell_h$ of part~(iii) is not treated). If moreover
$\eta\Delta R\ge\max\{\ell_h,\tfrac12(\ell_h+\ell_v)\}$, then $g$ is
nondecreasing and the set of amplification times is empty, or all of
$(0,\infty)$, or a half-line $(\tau_{\mathrm{amp}},\infty)$.
\end{enumerate}
\end{proposition}

\begin{proof}
(i) Put $z:=e^{\eta\lambda\Delta R}$. By
Theorem~\ref{thm:dynamics-tabular}(a),
$\pi^A_t(f)=p_f\beta_f^\lambda z/\bigl(p_h\beta_h^\lambda+(p_v\beta_v^\lambda+p_f\beta_f^\lambda)z\bigr)$
and $\pi^{\mathrm{no}}_t(f)=p_fz/(p_h+(p_v+p_f)z)$. Cross-multiplying and
cancelling the term $p_f\beta_f^\lambda z$ gives
\eqref{eq:equal-time-criterion}.

(ii) Divide \eqref{eq:equal-time-criterion} by $\beta_f^\lambda$ and write
$\varrho_v:=\beta_v/\beta_f\in[0,1)$, $\varrho_h:=\beta_h/\beta_f\ge1$:
the criterion is
$p_v(1-\varrho_v^\lambda)z>p_h(\varrho_h^\lambda-1)$. If
$\varrho_h=1$ the right-hand side vanishes and the left-hand side is
positive. If $\varrho_h>1$, the weighted arithmetic--geometric mean
inequality gives, for $\lambda\in(0,1)$,
$\varrho_v^\lambda<1-\lambda+\lambda\varrho_v$ and
$\varrho_h^\lambda<1-\lambda+\lambda\varrho_h$, i.e.\
$1-\varrho_v^\lambda>\lambda(1-\varrho_v)$ and
$\varrho_h^\lambda-1<\lambda(\varrho_h-1)$. Therefore the criterion holds
as soon as $p_v(1-\varrho_v)z\ge p_h(\varrho_h-1)$, i.e.\
$\eta\lambda\Delta R\ge T_{\mathrm{amp}}$. The two explicit cases follow
from $\lambda(t)=1-e^{-t/\eta}$.

(iii) and (iv). For $\beta_v>0$,
$g(\lambda)=\log(p_v/p_h)+\eta\Delta R\,\lambda+\log(1-e^{-\ell_v\lambda})-\log(e^{\ell_h\lambda}-1)$,
and $g(\lambda)\to\log\bigl(p_v\ell_v/(p_h\ell_h)\bigr)$ as
$\lambda\downarrow0$, which gives (iii) for $\beta_v>0$. For $\beta_v=0$
the left-hand side of the criterion tends to $p_v>0$ and the right-hand
side to $0$. Let $\varphi(y):=y/(e^y-1)$, $\varphi(0):=1$. Using
$\ell e^{\ell\lambda}/(e^{\ell\lambda}-1)=\ell+\varphi(\ell\lambda)/\lambda$,
\[
g'(\lambda)=\eta\Delta R-\ell_h+\frac{\varphi(\ell_v\lambda)-\varphi(\ell_h\lambda)}{\lambda}
=\eta\Delta R-\ell_h+\int_{\ell_h}^{\ell_v}\varphi'(c\lambda)\,dc .
\]
The function $\varphi$ is convex: for $y\neq0$,
$\varphi''(y)=e^y\,\mu(y)/(e^y-1)^3$ with $\mu(y)=(y-2)e^y+y+2$, and
$\mu(0)=\mu'(0)=0$, $\mu''(y)=ye^y$, so $\mu$ has the sign of $y$, as has
$(e^y-1)^3$. Hence $\lambda\mapsto\varphi'(c\lambda)$ is nondecreasing for
each $c>0$, and $g'$ is nondecreasing if $\ell_v\ge\ell_h$ and
nonincreasing if $\ell_v\le\ell_h$; the condition
$\ell_v\ge\ell_h$ is $\beta_f^2\ge\beta_v\beta_h$. For $\beta_v=0$,
$g'(\lambda)=\eta\Delta R-\ell_h-\varphi(\ell_h\lambda)/\lambda$, which is
increasing because $\varphi(\ell_h\lambda)/\lambda=\ell_h/(e^{\ell_h\lambda}-1)$
is decreasing; this is the convex case. The description of the amplification
set $\{g>0\}$ follows from convexity or concavity. Finally
$\varphi'(0)=-\tfrac12$ and $\varphi'(y)\to0$ as $y\to\infty$, so
$g'(0^+)=\eta\Delta R-\tfrac12(\ell_h+\ell_v)$ and
$g'(\infty)=\eta\Delta R-\ell_h$; $g'$ being monotone, it is bounded below
by the smaller of the two, which is nonnegative under the stated
condition (for $\beta_v=0$, $g'$ increases from $-\infty$, and the
condition cannot hold).
\end{proof}

\begin{remark}[No closed form, and re-entrance]\label{rem:reentrance}
Unlike the dominance times of Corollary~\ref{cor:hitting-times}, the
amplification time is the root of the transcendental equation $g=0$ and
has no closed form in general; part~(ii) shows that it is never later than
the time at which the effective pressure reaches the static threshold. In
the convex case the amplification set can be disconnected: for
$\beta=(\beta_h,\beta_v,\beta_f)=(1,\,5\cdot10^{-4},\,0.5)$,
$(p_h,p_v)=(0.8,\,0.084)$ and $\eta\Delta R=3$, a numerical evaluation gives
$g(0^+)\approx0.045$, $\min g\approx-0.037$ (at $\lambda\approx0.210$) and $g(1)\approx0.745$, with zeros at $\lambda\approx0.068$ and $0.360$ ($t/\eta\approx0.071$ and $0.446$), so that
operator~A amplifies stealth at very short times, then does not, then
amplifies it again. The effect is real but tiny: $\pi_A(f)-\pi_{\mathrm{no}}(f)$ is at most about $5\cdot10^{-5}$ at short times and about $-7.4\cdot10^{-4}$ in the trough.
\end{remark}

\paragraph{The regularised discard policy gradient.}
Part~(c2) of Theorem~\ref{thm:dynamics-tabular} has no anchor. With the
same anchor as in rule~(a), the rule becomes the natural gradient flow of
\[
V_{\mathrm{acc}}(\pi):=b(\pi)-\eta^{-1}\mathrm{KL}(\pi\,\|\,\pi_0),
\qquad\text{i.e.}\qquad
\dot\theta=a(\pi_\theta)-\eta^{-1}(\theta-\log p)\ \ (\mathrm{mod}\ \mathbf 1),
\quad\theta(0)=\log p .
\]
Its behaviour is qualitatively different from both rule~(a) and the
unregularised case.

\begin{proposition}[Regularised discard policy gradient]\label{prop:discard-kl}
Let $\pi_t$ be the flow above, let
$y(t):=\log\frac{\pi_t(f)}{\pi_t(h)}+L$ be the gain in log-odds of stealth
against honesty, and let $c:=(p_v+p_f)/p_h$.
\begin{enumerate}
\item[\textup{(i)}] \emph{Well-posedness.} The flow exists for all
$t\ge0$, $\bar Z(\pi_t)\ge\beta_fp_f/(p_f+p_v)$, and $\pi_t$ remains in a
compact subset of the open simplex. Moreover
$\pi_t(h)/p_h\le\pi_t(v)/p_v\le\pi_t(f)/p_f$ for all $t$.
\item[\textup{(ii)}] \emph{Equilibria.} $V_{\mathrm{acc}}(\pi_t)$ is
nondecreasing, and strictly increasing except at equilibria. The
equilibria are exactly the solutions in the open simplex of the
self-consistency equation
\begin{equation}\label{eq:self-consistent}
\pi(x)\;\propto\;p(x)\exp\bigl(\eta\,a_x(\pi)\bigr),
\end{equation}
equivalently the stationary points of $V_{\mathrm{acc}}$. The
$\omega$-limit set of $\pi_t$ is a nonempty connected set of equilibria; in
particular \eqref{eq:self-consistent} has a solution. Every solution
satisfies $\log\frac{\pi(f)}{\pi(h)}\ge-L+\eta\rho\Delta R$ and
$\log\frac{\pi(f)}{\pi(v)}=\log\frac{p_f}{p_v}+\eta(\beta_f-\beta_v)\beta_h\,\pi(h)\,\Delta R/\bar Z(\pi)^2$.
\item[\textup{(iii)}] \emph{Guaranteed dominance.}
$y(t)\ge\eta\lambda(t)\,\rho\,\Delta R$ for all $t$. Hence, if $L>0$ and
$\eta\rho\Delta R>L$, stealth dominates honesty from a time
$\tau\le-\eta\log\bigl(1-L/(\eta\rho\Delta R)\bigr)$ onwards.
\item[\textup{(iv)}] \emph{Trapping.} If there is $\bar y>0$ with
$\eta\rho\Delta R\,(1+c\,e^{\bar y})^2\le\bar y$, then $y(t)<\bar y$ for all
$t\ge0$. In particular, if such a $\bar y$ exists with $\bar y\le L$,
stealth never dominates honesty along the trajectory, nor at any of its
limit points.
\item[\textup{(v)}] \emph{Non-uniqueness and non-reachability.} Let
$p=(p_h,p_v,p_f)=(1-2\cdot10^{-4},10^{-4},10^{-4})$, $\beta_h=1$,
$\beta_f=0.1$, $\beta_v\in(0,0.1)$ and $\eta\Delta R=20$. Then the static
operator~A and the unmonitored static policy both have stealth dominating
honesty ($\eta\Delta R>L+\ell_h\approx11.51$), but \textup{(iv)} applies
with $\bar y=3$: along the regularised discard flow
$\log(\pi_t(f)/\pi_t(h))\le3-L<-6.2$ for all $t$. The global maximiser of
$V_{\mathrm{acc}}$ is a different equilibrium, which the flow started at
$\pi_0$ does not approach.
\end{enumerate}
\end{proposition}

\begin{proof}
Work in the gauge $\phi:=\theta-\log p$, $\dot\phi=a(\pi)-\phi/\eta$,
$\phi(0)=0$, which defines the same policy
$\pi_t(x)\propto p(x)e^{\phi_x(t)}$; $\phi$ is constant on strata. On the
maximal interval of existence,
$\phi_i(t)=\int_0^te^{-(t-s)/\eta}a_i(\pi_s)\,ds$.

(i) By \eqref{eq:discard-a} (which is an identity for $a$),
$a_h\le0\le a_v\le a_f$. Hence $\phi_h\le0\le\phi_v\le\phi_f$, which is the
chain of inequalities in the statement. It implies
$\pi_t(f)\ge\frac{p_f}{p_f+p_v}(\pi_t(f)+\pi_t(v))$, so
$\bar Z\ge\beta_h\pi_t(h)+\beta_f\pi_t(f)\ge\beta_f\bigl(\pi_t(h)+\tfrac{p_f}{p_f+p_v}(1-\pi_t(h))\bigr)\ge\beta_fp_f/(p_f+p_v)$.
Since $|\beta(x)(R(x)-b)|\le\beta_h\Delta R$, we get
$|a_x|\le M:=\beta_h\Delta R(p_f+p_v)/(\beta_fp_f)$ and $|\phi_x|\le\eta M$.
The solution is therefore global and
$\pi_t(x)\ge p(x)e^{-2\eta M}$.

(ii) By Lemma~\ref{lem:npg-softmax} with $V=V_{\mathrm{acc}}$ and
$\xi(\pi):=\nabla_\pi V_{\mathrm{acc}}=a(\pi)-\eta^{-1}\log(\pi/p)$ modulo
$\mathbf 1$, we have $\frac{d}{dt}\log\pi_t=\xi-\E_{\pi_t}[\xi]$ and
$\frac{d}{dt}V_{\mathrm{acc}}(\pi_t)=\sum_x\pi_t(x)(\xi_x-\E_{\pi_t}\xi)^2$.
Both vanish if and only if $\xi$ is constant, which is
\eqref{eq:self-consistent} and also the stationarity of
$V_{\mathrm{acc}}$ on the simplex. By (i) the trajectory is precompact in
the open simplex, where $V_{\mathrm{acc}}$ is continuous; by LaSalle's
invariance principle \cite{lasalle1976stability} its $\omega$-limit set is
nonempty, compact, connected, invariant and contained in the set where
$\frac{d}{dt}V_{\mathrm{acc}}=0$. At a solution of
\eqref{eq:self-consistent},
$\log(\pi(f)/\pi(h))=-L+\eta(a_f-a_h)$ and
$\log(\pi(f)/\pi(v))=\log(p_f/p_v)+\eta(a_f-a_v)$; the bound
$a_f-a_h\ge\rho\Delta R$ and the value of $a_f-a_v$ were obtained in the
proof of Theorem~\ref{thm:dynamics-tabular}(c2).

(iii) $y=\phi_f-\phi_h=\int_0^te^{-(t-s)/\eta}(a_f-a_h)\,ds\ge\rho\Delta R\,\eta(1-e^{-t/\eta})$,
and $\log(\pi_t(f)/\pi_t(h))=y(t)-L$.

(iv) Let $\mu:=m/(\beta_h\pi(h))$ with $m=\beta_v\pi(v)+\beta_f\pi(f)$.
From \eqref{eq:discard-a},
$a_f-a_h=\beta_h\Delta R(\beta_f\pi(h)+m)/\bar Z^2
=\Delta R\,(\rho+\mu)/\bigl(\pi(h)(1+\mu)^2\bigr)\le\Delta R(\rho+\mu)/\pi(h)$.
Along the trajectory, (i) gives
$\pi(v)/\pi(h)\le(p_v/p_h)e^{y}$ and $\pi(f)/\pi(h)=(p_f/p_h)e^{y}$, so
$\mu\le\rho\,c\,e^{y}$ (using $\beta_v\le\beta_f$) and
$1/\pi(h)\le1+c\,e^{y}$. Therefore
$\dot y=(a_f-a_h)-y/\eta\le\Gamma(y):=\rho\Delta R(1+c\,e^{y})^2-y/\eta$.
Since $\Gamma(0)>0\ge\Gamma(\bar y)$, $\Gamma$ has a smallest zero
$z^\circ\in(0,\bar y]$. The solution $z$ of $\dot z=\Gamma(z)$, $z(0)=0$,
stays below the constant solution $z^\circ$ by uniqueness, and the
comparison lemma for scalar differential inequalities gives
$y(t)\le z(t)<z^\circ\le\bar y$.

(v) Here $L=\log(p_h/p_f)\approx9.2101$, $\ell_h=\log10$, $\rho=0.1$,
$c\approx2.0004\cdot10^{-4}$, and
$\eta\rho\Delta R(1+c\,e^{3})^2\approx2.016\le3$, so (iv) applies with
$\bar y=3$. Let $\mathcal T$ be the set of policies with
$\pi(f)/\pi(h)\le(p_f/p_h)e^{3}$ and $\pi(v)/\pi(h)\le(p_v/p_h)e^{3}$; by
(i) and (iv) the trajectory and its limit points lie in $\mathcal T$. On
$\mathcal T$, $b=\Delta R\,\mu/(1+\mu)\le\Delta R\,\rho\,c\,e^{3}$, so
$\eta V_{\mathrm{acc}}\le\eta b\le20\cdot0.1\cdot c\,e^{3}<0.01$. On the
other hand, for the policy $\pi^\circ:=p(\cdot\mid\mathcal X_f)$ (the
pretrained policy conditioned on the stealth stratum) we get
$b(\pi^\circ)=\Delta R$ and $\mathrm{KL}(\pi^\circ\|p)=\log(1/p_f)$, so
$\eta V_{\mathrm{acc}}(\pi^\circ)=20-\log(10^{4})>10$. Since $\beta_v>0$,
$\bar Z\ge\beta_v$ on the whole simplex, $b$ is $C^1$ on a neighbourhood of
the closed simplex, and $V_{\mathrm{acc}}$ is continuous there; it attains
its maximum at some $\pi^\star$. This maximiser is interior: if
$\pi^\star(x)=0$ for some $x$, then along
$\pi_\varepsilon:=(1-\varepsilon)\pi^\star+\varepsilon p$ the function $b$
changes by $O(\varepsilon)$ while
$\frac{d}{d\varepsilon}\mathrm{KL}(\pi_\varepsilon\|p)
=\sum_{x'}(p(x')-\pi^\star(x'))\log\frac{\pi_\varepsilon(x')}{p(x')}\to-\infty$
as $\varepsilon\downarrow0$, so $V_{\mathrm{acc}}(\pi_\varepsilon)>V_{\mathrm{acc}}(\pi^\star)$
for small $\varepsilon$, a contradiction. Hence $\pi^\star$ is a stationary
point, i.e.\ an equilibrium, with
$\eta V_{\mathrm{acc}}(\pi^\star)>10$, so $\pi^\star\notin\mathcal T$,
whereas all limit points of the trajectory are equilibria in $\mathcal T$.
The static claims are $\eta\Delta R=20>L+\ell_h$ and $20>L$.
\end{proof}

\begin{remark}[What Theorem~\ref{thm:dynamics-tabular} says about the filter]
\label{rem:taxonomy}
Rules (a)--(c) are three different dynamical systems; they do not even
share their limits. Several of them produce policies of the form
$\pi\propto p\,\beta^{\,\mathsf a}\,e^{\mathsf bR}$, for which
$\log(\pi(f)/\pi(h))=-L+\mathsf b\,\Delta R-\mathsf a\,\ell_h$ and
$\log(\pi(f)/\pi(v))=\log(p_f/p_v)+\mathsf a\,\ell_v$. The pair
$(\mathsf a,\mathsf b)$ (filter exponent, pressure) is
$(1,\eta)$ for the static operator~A;
$(\lambda(t),\eta\lambda(t))$ for rule~(a) applied to operator~A;
$(k,k\gamma)$ for rule~(b) with $\omega(R)=e^{\gamma R}$; and
$(c_k,c_k\gamma)$ for its anchored variant. Along the static operator
$\mathsf a=1$ is fixed, so the filter term $\ell_h$ is eventually
overwhelmed as $\eta$ grows. Along rule~(b) the ratio
$\mathsf b/\mathsf a=\gamma$ is fixed: the filter \emph{compounds} at the
same rate as the reward pressure, the sign of
$\vartheta=\gamma\Delta R-\ell_h$ decides the outcome whatever the number
of rounds, and the odds of stealth against visible cheating grow by
$\beta_f/\beta_v$ at \emph{every} round. If $\vartheta<0$ the compounding
works against stealth relative to honesty.

Rule~(c), in its natural-gradient form, is of a different nature. By (c1) it ascends the accepted-sample
mean reward, in which the filter multiplies the advantage of each stratum
instead of tilting the target distribution. Without anchor, the growth rate of the odds of stealth against honesty is at least $\rho=\beta_f/\beta_h$ times $\Delta R$; the filter never reverses this growth and may even accelerate it (at $\pi=(0.01,0,0.99)$ with $\beta=(1,0,0.1)$, $a_f-a_h=8.4>\Delta R=1$); and it shifts mass from
visible to stealthy cheating only by a bounded amount, at a rate
proportional to the remaining honest mass
(Theorem~\ref{thm:dynamics-tabular}(c2)). With the anchor, however, the
statement ``stealth wins for every $\beta_f>0$'' is false
(Proposition~\ref{prop:discard-kl}): heuristically, the filter then rescales the
\emph{pressure}, dominance from the pretrained policy is guaranteed when
$\eta\rho\Delta R>L$ and excluded when
$\eta\rho\Delta R(1+ce^{\bar y})^2\le\bar y$ for some $\bar y\le L$, and
several equilibria can coexist. The two thresholds are far apart: in the instance of Proposition~\ref{prop:discard-kl}(v) with $\beta_v=0.05$, the trapping bound is nearly tight (largest $y$ equal to $2.003$ against $\bar y=2.006$), but dominance from the pretrained policy appears for $\eta\rho\Delta R$ between $5$ and $6$, well below the guaranteed threshold $L=9.21$. The global maximiser in that instance is close to $(0,0.5,0.5)$, so it does not enrich stealth among cheaters. The identity
$a_f-a_h=\bigl(\beta_f\Delta R+(\beta_h-\beta_f)b\bigr)/\bar Z$ shows where
the multiplicity comes from: the first term is, up to the factor $1/\bar Z$, what operator~B with $C=0$
would give, and the second is a positive feedback, absent from
operator~B, which grows with the accepted-sample mean reward $b$ because
intercepted samples are dropped rather than scored.

For selection rules whose weights depend on the current policy, part~(b)
extends with round-dependent increments. For best-of-$N$ selection (at each
round, $N$ accepted samples are drawn, one of highest reward is kept, ties
being broken uniformly, and $\pi_{k+1}$ is the law of the kept sample), the
kept sample has the accepted law reweighted by
$\omega_k(\Delta R)=(1-(1-\chi_k)^N)/\chi_k$ and
$\omega_k(0)=(1-\chi_k)^{N-1}$, where $\chi_k\in(0,1)$ is the
accepted-sample share of cheating at round $k$. The increment of
$\log(\pi_k(f)/\pi_k(h))$ at round $k$ is therefore
$\log\sum_{i=0}^{N-1}(1-\chi_k)^{-i}-\ell_h\ge\log N-\ell_h$, and
stealth eventually dominates whenever $N>\beta_h/\beta_f$. We do not
characterise the case $N\le\beta_h/\beta_f$, in which the outcome depends
on the initial policy.
\end{remark}

% ----------------------------------------------------------------------
\subsection{Two-channel dynamics: action and chain of thought}
\label{sec:two-channel}

We now let the policy act through two channels with separate parameters.
The \emph{action channel} decides whether to hack, with probability
$q=\sigma(u)$; the \emph{chain-of-thought channel} decides whether the
reasoning is obfuscated, with probability $s=\sigma(w)$, independently of
the action; here $\sigma(z)=1/(1+e^{-z})$. The pretrained model has
$q_0=\sigma(u_0)$ and $s_0=\sigma(w_0)$, both in $(0,1)$. The monitor reads
the chain of thought: a hack survives with probability $\beta_v$ if the
reasoning is transparent and $\beta_f$ if it is obfuscated; an honest
trajectory survives with probability $\beta_h$ in both cases. Obfuscation
costs $\kappa\ge0$ whatever the action. The correspondence with the strata
is: honest $=h$, transparent hack $=v$, obfuscated hack $=f$.

Under operator~B the effective reward is $-(1-\beta_h)C$ for an honest
trajectory, $\beta_v\Delta R-(1-\beta_v)C$ for a transparent hack and
$\beta_f\Delta R-(1-\beta_f)C$ for an obfuscated hack, minus $\kappa$ when
the reasoning is obfuscated. Up to the additive constant $-(1-\beta_h)C$,
the regularised objective restricted to product policies is therefore
\begin{equation}\label{eq:two-channel-J}
\begin{gathered}
J(q,s)=q\,(A_0+Ds)-\kappa s-\eta^{-1}\bigl[\mathrm{kl}(q\|q_0)+\mathrm{kl}(s\|s_0)\bigr],\\
A_0:=\beta_v\Delta R-(\beta_h-\beta_v)C,\qquad
D:=(\beta_f-\beta_v)(\Delta R+C)>0,
\end{gathered}
\end{equation}
where $\mathrm{kl}(a\|b)=a\log\frac ab+(1-a)\log\frac{1-a}{1-b}$ (for
$\beta_h=1$, $A_0=\beta_v\Delta R-(1-\beta_v)C$). The Fisher information of
a Bernoulli variable with logit $u$ is $\sigma'(u)$, and
$\partial_qJ=A_0+Ds-\eta^{-1}(u-u_0)$,
$\partial_sJ=Dq-\kappa-\eta^{-1}(w-w_0)$, so the natural gradient flow of
$J$ in the coordinates $x=(u,w)$ is
\begin{equation}\label{eq:two-channel-flow}
\dot u=f_1(u,w):=A_0+D\,\sigma(w)-\frac{u-u_0}{\eta},\qquad
\dot w=f_2(u,w):=D\,\sigma(u)-\kappa-\frac{w-w_0}{\eta},
\end{equation}
started at $x_0=(u_0,w_0)$ unless stated otherwise. The vector field
$f=(f_1,f_2)$ is \emph{cooperative}: $\partial_wf_1=D\sigma'(w)>0$ and
$\partial_uf_2=D\sigma'(u)>0$. Inequalities between points of
$\mathbb R^2$ are meant coordinatewise. Define the increasing maps
\[
U(w):=u_0+\eta\bigl(A_0+D\sigma(w)\bigr),\qquad
W(u):=w_0+\eta\bigl(D\sigma(u)-\kappa\bigr),\qquad
\Phi:=W\circ U .
\]

\begin{theorem}[Two-channel dynamics under operator~B]
\label{thm:dynamics-two-channel}
Consider the flow \eqref{eq:two-channel-flow} with $\eta>0$, $D>0$,
$\kappa\ge0$ and $A_0\in\mathbb R$.

\smallskip\noindent\textup{(a) Global behaviour.}
Solutions exist for all $t\ge0$ from every initial condition and are
bounded; the rectangle
$[u_0+\eta A_0,\,u_0+\eta(A_0+D)]\times[w_0-\eta\kappa,\,w_0+\eta(D-\kappa)]$
is forward invariant and attracts every solution. The function
$\hat J(u,w):=J(\sigma(u),\sigma(w))$ satisfies
$\frac{d}{dt}\hat J=\sigma'(u)f_1^2+\sigma'(w)f_2^2\ge0$, with equality
only at equilibria. Every solution converges to an equilibrium.

\smallskip\noindent\textup{(b) Equilibria.}
\begin{enumerate}
\item[\textup{(b1)}] The map $(u,w)\mapsto w$ is a bijection from the set
of equilibria onto the set of fixed points of $\Phi$, with inverse
$w\mapsto(U(w),w)$. Equilibria are totally ordered: if
$(u,w)\ne(u',w')$ are equilibria then either $u<u'$ and $w<w'$, or $u>u'$
and $w>w'$.
\item[\textup{(b2)}] $\Phi'$ is strictly increasing on $(-\infty,w^\dagger]$
and strictly decreasing on $[w^\dagger,\infty)$ for some
$w^\dagger\in\mathbb R$. Consequently there are at least one and at most
three equilibria.
\item[\textup{(b3)}] The Jacobian of $f$ at an equilibrium $(U(w^\ast),w^\ast)$
has the real eigenvalues $-\bigl(1\mp\sqrt{\Phi'(w^\ast)}\bigr)/\eta$. The
equilibrium is a hyperbolic stable node if $\Phi'(w^\ast)<1$, a hyperbolic
saddle if $\Phi'(w^\ast)>1$, and non-hyperbolic if $\Phi'(w^\ast)=1$. If
there are exactly three equilibria $x_1<x_2<x_3$, then $x_1$ and $x_3$ are
hyperbolic stable nodes and $x_2$ is a hyperbolic saddle.
\item[\textup{(b4)}] If $\eta D<4$, the equilibrium $x^\ast$ is unique and
every solution satisfies
\[
\|x(t)-x^\ast\|_\infty\le e^{-(1/\eta-D/4)\,t}\,\|x(0)-x^\ast\|_\infty .
\]
\end{enumerate}

\smallskip\noindent\textup{(c) Order certificates.}
Let $x(t)$ be the solution started at $x_0$.
\begin{enumerate}
\item[\textup{(c1)}] If $\bar x\ge x_0$ and $f(\bar x)\le0$, then
$x(t)\le\bar x$ for all $t\ge0$. If $\underline x\le x_0$ and
$f(\underline x)\ge0$, then $x(t)\ge\underline x$ for all $t\ge0$.
\item[\textup{(c2)}] If $f(x_0)\le0$, then $x(t)$ is nonincreasing in both
coordinates and converges to the largest equilibrium $\le x_0$.
\item[\textup{(c3)}] If $f(x_0)\ge0$, then $x(t)$ is nondecreasing in both
coordinates and converges to the smallest equilibrium $\ge x_0$.
\end{enumerate}

\smallskip\noindent\textup{(d) Local rates.}
Let $x^\ast=(U(w^\ast),w^\ast)$ be an equilibrium with $\Phi'(w^\ast)<1$ and
let $\mu^\ast:=\bigl(1-\sqrt{\Phi'(w^\ast)}\bigr)/\eta$. For every
$\mu<\mu^\ast$ there are a neighbourhood $\mathcal N$ of $x^\ast$ and a
constant $K$ such that $\|x(t)-x^\ast\|\le Ke^{-\mu t}\|x(0)-x^\ast\|$ for
every solution started in $\mathcal N$. Here
$\Phi'(w^\ast)=\eta^2D^2\,q^\ast(1-q^\ast)\,s^\ast(1-s^\ast)$ with
$q^\ast=\sigma(U(w^\ast))$, $s^\ast=\sigma(w^\ast)$.
\end{theorem}

\begin{proof}
\emph{(b1).} $f(u,w)=0$ if and only if $u=U(w)$ and $w=W(u)$, i.e.\ $w$ is a
fixed point of $\Phi$ and $u=U(w)$. Since $U$ is strictly increasing, two
distinct equilibria have distinct $w$-coordinates and their
$u$-coordinates are ordered in the same way.

\emph{(b2).} We have
$\Phi'(w)=\eta D\,\sigma'(U(w))\cdot\eta D\,\sigma'(w)>0$. Using
$(\log\sigma')'=1-2\sigma$ and $U'(w)=\eta D\sigma'(w)$,
\[
\frac{d}{dw}\log\Phi'(w)=\bigl(1-2\sigma(U(w))\bigr)\,\eta D\,\sigma'(w)+1-2\sigma(w)
=s(1-s)\Bigl[\psi(s)-\eta D\tanh\!\Bigl(\tfrac{u_0+\eta A_0+\eta Ds}{2}\Bigr)\Bigr],
\]
where $s=\sigma(w)\in(0,1)$, $\psi(s):=\frac{1-2s}{s(1-s)}=\frac1s-\frac1{1-s}$,
and we used $2\sigma(z)-1=\tanh(z/2)$. The function $\psi$ is strictly
decreasing from $+\infty$ to $-\infty$ on $(0,1)$, while
$s\mapsto\eta D\tanh((u_0+\eta A_0+\eta Ds)/2)$ is strictly increasing and
bounded. Hence the bracket has exactly one zero $s^\dagger$, is positive
for $s<s^\dagger$ and negative for $s>s^\dagger$; this is the claim with
$w^\dagger=\sigma^{-1}(s^\dagger)$. It follows that $\Phi'-1$ has at most
two zeros, and by Rolle's theorem $\Phi-\mathrm{id}$ has at most three.
It has at least one because $\Phi$ is continuous and bounded, so that
$\Phi(w)-w\to\pm\infty$ as $w\to\mp\infty$.

\emph{(b3).} At an equilibrium $(u^\ast,w^\ast)$ the Jacobian of $f$ is
\[
\begin{pmatrix}-1/\eta& D\sigma'(w^\ast)\\ D\sigma'(u^\ast)&-1/\eta\end{pmatrix},
\quad\text{with eigenvalues}\quad
-\frac1\eta\pm D\sqrt{\sigma'(u^\ast)\sigma'(w^\ast)}
=-\frac{1\mp\sqrt{\Phi'(w^\ast)}}{\eta},
\]
since $\Phi'(w^\ast)=\eta^2D^2\sigma'(u^\ast)\sigma'(w^\ast)$ by the
formula in the proof of (b2). The classification follows. If
$\Phi-\mathrm{id}$ has three zeros $w_1<w_2<w_3$, Rolle's theorem gives
zeros $\xi_1\in(w_1,w_2)$ and $\xi_2\in(w_2,w_3)$ of $\Phi'-1$, which by
(b2) are its only zeros, and $\Phi'-1>0$ exactly on $(\xi_1,\xi_2)$. Hence
$\Phi'(w_1)<1$, $\Phi'(w_2)>1$, $\Phi'(w_3)<1$.

\emph{(a).} The field $f$ is globally Lipschitz, so solutions are global.
Variation of constants gives
\[
u(t)=u_0+e^{-t/\eta}(u(0)-u_0)+\int_0^te^{-(t-\varsigma)/\eta}\bigl(A_0+D\sigma(w(\varsigma))\bigr)\,d\varsigma ,
\]
and the integral lies in $\eta(1-e^{-t/\eta})[A_0,A_0+D]$; similarly for
$w$ with $[-\kappa,D-\kappa]$. This proves boundedness and the statement
on the rectangle. Next, $\partial_u\hat J=\sigma'(u)\,\partial_qJ=\sigma'(u)f_1$
and $\partial_w\hat J=\sigma'(w)f_2$ by the formulas preceding
\eqref{eq:two-channel-flow}, so
$\frac{d}{dt}\hat J=\sigma'(u)f_1^2+\sigma'(w)f_2^2$, which vanishes only
where $f=0$. By LaSalle's invariance principle
\cite{lasalle1976stability}, the $\omega$-limit set of a (bounded) solution
is nonempty, connected and made of equilibria; by (b2) there are at most
three equilibria, so it is a single point.

\emph{(b4).} If $\eta D<4$ then $\Phi'\le(\eta D)^2/16<1$, so
$\Phi-\mathrm{id}$ is strictly decreasing and the equilibrium is unique.
Let $x,\tilde x$ be two solutions and $\delta:=x-\tilde x$. By the mean
value theorem $\dot\delta=-\delta/\eta+N(t)\delta$, where $N(t)$ has zero
diagonal and off-diagonal entries in $[0,D/4]$ (as $0<\sigma'\le\frac14$).
For $z:=e^{t/\eta}\delta$ we get
$\|z(t)\|_\infty\le\|z(0)\|_\infty+\frac D4\int_0^t\|z\|_\infty$, and
Gr\"onwall's lemma yields
$\|\delta(t)\|_\infty\le e^{-(1/\eta-D/4)t}\|\delta(0)\|_\infty$. Take
$\tilde x\equiv x^\ast$.

\emph{(c1).} Let $\bar x=(\bar u,\bar w)\ge x_0$ with $f(\bar x)\le0$. For
$\varepsilon>0$ let $x^\varepsilon$ solve
$\dot x=f(x)-\varepsilon(1,1)$, $x^\varepsilon(0)=x_0$, and let
$t_1:=\sup\{t\ge0:\ x^\varepsilon\le\bar x\text{ on }[0,t]\}$. Suppose
$t_1<\infty$. Then $x^\varepsilon(t_1)\le\bar x$. For a coordinate in which
equality holds, say $u^\varepsilon(t_1)=\bar u$, cooperativity gives
$\dot u^\varepsilon(t_1)=f_1(\bar u,w^\varepsilon(t_1))-\varepsilon\le f_1(\bar u,\bar w)-\varepsilon<0$,
so $u^\varepsilon<\bar u$ on some interval $(t_1,t_1+\delta]$; a coordinate
in which the inequality is strict remains strictly below on such an
interval by continuity. Hence $x^\varepsilon\le\bar x$ on
$[0,t_1+\delta]$, contradicting the definition of $t_1$. Thus
$x^\varepsilon(t)\le\bar x$ for all $t$, and letting $\varepsilon\downarrow0$
(continuous dependence on parameters, on compact time intervals) gives
$x(t)\le\bar x$. The second statement is symmetric.

\emph{(c2) and (c3).} Let $y:=\dot x=f(x(t))$. Then
$\dot y=Df(x(t))\,y=-y/\eta+N(t)y$ with
$N(t)=\begin{pmatrix}0&D\sigma'(w(t))\\D\sigma'(u(t))&0\end{pmatrix}$
entrywise nonnegative. For $z:=e^{t/\eta}y$ we have $\dot z=N(t)z$, and the
Picard iterates $z^{(0)}\equiv z(0)$,
$z^{(n+1)}(t)=z(0)+\int_0^tN(\varsigma)z^{(n)}(\varsigma)\,d\varsigma$
converge to $z$. If $z(0)=f(x_0)\ge0$ all iterates are $\ge0$, hence
$\dot x(t)\ge0$ for all $t$; if $f(x_0)\le0$, apply this to $-z$. In both
cases $x(t)$ is monotone and bounded, so it converges, necessarily to an
equilibrium $x^\ast$, with $x^\ast\ge x_0$ in case (c3) and $x^\ast\le x_0$
in case (c2). In case (c3), any equilibrium $e\ge x_0$ satisfies
$f(e)=0\le0$, so $x(t)\le e$ by (c1) and $x^\ast\le e$. Case (c2) is
symmetric.

\emph{(d).} By (b3) the Jacobian at $x^\ast$ has the two real, distinct,
negative eigenvalues $-(1\mp\sqrt{\Phi'(w^\ast)})/\eta$, the larger being
$-\mu^\ast$; the statement is the standard linearisation estimate for
hyperbolic sinks. The formula for $\Phi'(w^\ast)$ is the one in the proof
of (b2).
\end{proof}

\paragraph{Explicit certificates and the penalty window.}
Let $\bar\beta_0:=\beta_v+s_0(\beta_f-\beta_v)$ be the survival probability
of a hack under the pretrained chain-of-thought channel; note
$\bar\beta_0<\beta_f\le\beta_h$. Define
\[
C_{\mathrm{low}}:=\frac{\bar\beta_0\,\Delta R}{\beta_h-\bar\beta_0},\qquad
C_{\mathrm{high}}:=\frac{\kappa}{q_0(\beta_f-\beta_v)}-\Delta R,\qquad
\kappa^\dagger:=\frac{q_0(\beta_f-\beta_v)\,\beta_h\,\Delta R}{\beta_h-\bar\beta_0}.
\]
We write $z^+:=\max(z,0)$.

\begin{corollary}[Certificates from the pretrained model]
\label{cor:penalty-window}
Let $(q(t),s(t))=(\sigma(u(t)),\sigma(w(t)))$ be the solution of
\eqref{eq:two-channel-flow} started at $x_0$.
\begin{enumerate}
\item[\textup{(i)}] Let $\bar u:=u_0+\eta(A_0+Ds_0)^+$. If
$D\,\sigma(\bar u)\le\kappa$, then $q(t)\le\sigma(\bar u)$ and
$s(t)\le s_0$ for all $t\ge0$.
\item[\textup{(ii)}] Let $\bar w:=w_0+\eta(Dq_0-\kappa)^+$. If
$A_0+D\,\sigma(\bar w)\le0$, then $q(t)\le q_0$ and
$s(t)\le\sigma(\bar w)$ for all $t\ge0$.
\item[\textup{(iii)}] \emph{Safe window.} The following are equivalent:
\textup{(1)} $q(t)\le q_0$ and $s(t)\le s_0$ for all $t\ge0$;
\textup{(2)} $A_0+Ds_0\le0$ and $Dq_0\le\kappa$;
\textup{(3)} $C_{\mathrm{low}}\le C\le C_{\mathrm{high}}$. When they hold,
$q(t)$ and $s(t)$ are nonincreasing.
\item[\textup{(iv)}] \emph{Reverse window.} If
$C_{\mathrm{high}}\le C\le C_{\mathrm{low}}$, then $q(t)$ and $s(t)$ are
nondecreasing, and the limit equilibrium $(u^\ast,w^\ast)$ satisfies
$u^\ast\ge u_0+\eta(A_0+Ds_0)$ and
\[
w^\ast\;\ge\;w_0+\eta\bigl(q_0(\beta_f-\beta_v)(\Delta R+C)-\kappa\bigr),
\]
a lower bound that increases linearly with the penalty $C$.
\item[\textup{(v)}] \emph{Dichotomy.} The safe window is nonempty if and
only if $\kappa\ge\kappa^\dagger$; the reverse window
$[C_{\mathrm{high}}^+,C_{\mathrm{low}}]$ is nonempty if and only if
$\kappa\le\kappa^\dagger$. For every state $(u,w)$,
$\partial f_1/\partial C=-(\beta_h-\beta_v)+(\beta_f-\beta_v)\sigma(w)<0$ and
$\partial f_2/\partial C=(\beta_f-\beta_v)\sigma(u)>0$.
\end{enumerate}
\end{corollary}

\begin{proof}
(i) Take $\bar x=(\bar u,w_0)\ge x_0$. Then
$f_1(\bar x)=(A_0+Ds_0)-(A_0+Ds_0)^+\le0$ and
$f_2(\bar x)=D\sigma(\bar u)-\kappa\le0$; apply
Theorem~\ref{thm:dynamics-two-channel}(c1). (ii) Take $\bar x=(u_0,\bar w)$:
$f_1(\bar x)=A_0+D\sigma(\bar w)\le0$ and
$f_2(\bar x)=(Dq_0-\kappa)-(Dq_0-\kappa)^+\le0$. (iii) Condition (2) is
$f(x_0)\le0$, which implies (1) and the monotonicity by
Theorem~\ref{thm:dynamics-two-channel}(c2). Conversely, if $f_1(x_0)>0$
then $\dot u(0)>0$ and $q(t)>q_0$ for small $t>0$, and likewise for $w$; so
(1) implies (2). Finally
$A_0+Ds_0=\bar\beta_0\Delta R-(\beta_h-\bar\beta_0)C$, which is $\le0$ if
and only if $C\ge C_{\mathrm{low}}$, and $Dq_0\le\kappa$ if and only if
$\Delta R+C\le\kappa/(q_0(\beta_f-\beta_v))$. (iv) The hypothesis is
$f(x_0)\ge0$; apply Theorem~\ref{thm:dynamics-two-channel}(c3). At the
limit equilibrium $u^\ast=U(w^\ast)\ge U(w_0)$ and
$w^\ast=W(u^\ast)\ge W(u_0)=w_0+\eta(Dq_0-\kappa)$. (v)
$C_{\mathrm{low}}\le C_{\mathrm{high}}$ if and only if
$\Delta R+C_{\mathrm{low}}=\beta_h\Delta R/(\beta_h-\bar\beta_0)\le\kappa/(q_0(\beta_f-\beta_v))$,
i.e.\ $\kappa\ge\kappa^\dagger$; note $C_{\mathrm{low}}>0$. The derivatives
are computed from \eqref{eq:two-channel-J}, and the first is negative
because $\beta_v+(\beta_f-\beta_v)\sigma(w)<\beta_f\le\beta_h$.
\end{proof}

\begin{remark}[The rejection--penalty dilemma in dynamic form]
\label{rem:dilemma}
Corollary~\ref{cor:penalty-window}(v) is the differential form of the
dilemma: raising the penalty always lowers the instantaneous drive towards
hacking and always raises the instantaneous drive towards obfuscation, the
latter in proportion to the current hack rate, because the penalty
increases $D$, the payoff of obfuscating conditional on hacking. Its
integrated consequences are the two windows. If obfuscation is costly
enough ($\kappa\ge\kappa^\dagger$), penalties in
$[C_{\mathrm{low}},C_{\mathrm{high}}]$ keep both the hack rate and the
obfuscation rate below their pretrained values for all time, and by
(iii) this property fails for every penalty outside the window. If
$\kappa<\kappa^\dagger$, \emph{no} penalty has this property, and for
penalties in the reverse window both rates grow monotonically, the
obfuscation logit by at least an amount that increases linearly in $C$.

We stress what is \emph{not} proved. For $C>C_{\mathrm{high}}$ (with
$C\ge C_{\mathrm{low}}$) the obfuscation rate does exceed $s_0$ at small
positive times, but this need not lead to the hack--obfuscated equilibrium:
the same penalty also depresses the hack rate, and the excursion can be
negligible; part~(ii) gives a quantitative bound in a range of such
penalties. In the instance of Proposition~\ref{prop:nonreach}(iii) below
($C_{\mathrm{high}}=9$), a numerical integration with $C=100$ gives
$\sup_ts(t)\approx0.01007$ for $s_0=0.01$. The reverse window, where hacking is profitable on average under the pretrained model, is the regime in which the growth of obfuscation with the penalty is \emph{proved} monotone. Outside both windows a larger penalty can still move the pretrained model from the honest equilibrium to the hack--obfuscated one: with $q_0=s_0=0.8$, $\beta=(1,0.04,0.85)$, $\Delta R=1$, $\kappa=1.2$ and $\eta=40$ (so $C_{\mathrm{high}}=0.852<C_{\mathrm{low}}=2.205$), the limit from the pretrained model is honest, $(q,s)\approx(4\cdot10^{-6},0)$, at $C=0.40$, and hack--obfuscated at $C=0.45$, $0.50$ and $0.60$ ($s\approx0.59$, $0.88$ and $0.995$), a fold occurring between $0.40$ and $0.45$. Such flips were found in $13$ of $800$ random draws over a wide parameter range.
\end{remark}

\begin{proposition}[Bistability]\label{prop:bistability}
Suppose $0<\kappa<D$ and that the two inequalities
\begin{align}
\sigma\Bigl(u_0+\eta\bigl(A_0+D\,\sigma(w_0-\tfrac12\eta\kappa)\bigr)\Bigr)&<\frac{\kappa}{2D},
\label{eq:cond-L}\\
\sigma\Bigl(u_0+\eta\bigl(A_0+D\,\sigma(w_0+\tfrac12\eta(D-\kappa))\bigr)\Bigr)&>\frac{D+\kappa}{2D}
\label{eq:cond-H}
\end{align}
hold. Then there are exactly three equilibria $x_1<x_2<x_3$, with
$w_1<w_0-\frac12\eta\kappa<w_2<w_0+\frac12\eta(D-\kappa)<w_3$; $x_1$ and
$x_3$ are hyperbolic stable nodes and $x_2$ is a hyperbolic saddle.
Moreover, if $A_0<0<A_0+D$ and $0<\kappa<D$, then \eqref{eq:cond-L} and
\eqref{eq:cond-H} hold for all sufficiently large $\eta$. In terms of the
primitives, $A_0<0<A_0+D$ reads
$\beta_v\Delta R<(\beta_h-\beta_v)C$ and
$\beta_f\Delta R>(\beta_h-\beta_f)C$; when $\beta_f=\beta_h$ the second
inequality always holds, and the bistable region contains
$\{(\eta,C):\ C>\beta_v\Delta R/(\beta_h-\beta_v),\ \eta\ge\eta^\ast(C)\}$
for some finite $\eta^\ast(C)$, provided $0<\kappa<(\beta_f-\beta_v)\Delta R$.
\end{proposition}

\begin{proof}
Let $w_-:=w_0-\frac12\eta\kappa$ and $w_+:=w_0+\frac12\eta(D-\kappa)$, so
$w_-<w_+$. Then
$\Phi(w_-)-w_-=\eta D\bigl[\sigma(U(w_-))-\kappa/(2D)\bigr]<0$ by
\eqref{eq:cond-L} and
$\Phi(w_+)-w_+=\eta D\bigl[\sigma(U(w_+))-(D+\kappa)/(2D)\bigr]>0$ by
\eqref{eq:cond-H}. Since $\Phi(w)-w\to\pm\infty$ as $w\to\mp\infty$, the
intermediate value theorem gives a zero in each of $(-\infty,w_-)$,
$(w_-,w_+)$ and $(w_+,\infty)$, and
Theorem~\ref{thm:dynamics-two-channel}(b2)--(b3) gives the rest. For the
asymptotic statement, as $\eta\to\infty$:
$\eta D\sigma(w_0-\frac12\eta\kappa)\le\eta De^{w_0-\eta\kappa/2}\to0$ and
$\eta A_0\to-\infty$, so the left-hand side of \eqref{eq:cond-L} tends to
$0<\kappa/(2D)$; and $\sigma(w_0+\frac12\eta(D-\kappa))\to1$, so
$A_0+D\sigma(\cdot)\to A_0+D>0$, the argument of the outer $\sigma$ in
\eqref{eq:cond-H} tends to $+\infty$ and the left-hand side to
$1>(D+\kappa)/(2D)$. The reformulations follow from
\eqref{eq:two-channel-J}: $A_0+D=\beta_f\Delta R-(\beta_h-\beta_f)C$, and
$D\ge(\beta_f-\beta_v)\Delta R$.
\end{proof}

\begin{proposition}[Non-reachability of the global maximiser]
\label{prop:nonreach}
Suppose
\begin{equation}\label{eq:nonreach-cond}
A_0+Ds_0\le0,\qquad Dq_0\le\kappa,\qquad
A_0+D-\kappa>\eta^{-1}\log\frac{1}{q_0s_0}.
\end{equation}
\begin{enumerate}
\item[\textup{(i)}] The solution started at the pretrained model satisfies
$q(t)\le q_0$, $s(t)\le s_0$ for all $t$ and converges to an equilibrium
$x_\star\le x_0$ with $J(x_\star)\le0$.
\item[\textup{(ii)}] $J$ attains its maximum over $[0,1]^2$ at an interior
point, which is an equilibrium $x^\star$ with $J(x^\star)>0$; and
$x^\star>x_\star$ in both coordinates. Thus the global maximiser of the regularised
objective is an equilibrium with strictly larger hack and obfuscation
rates than $x_\star$, and it is not reached from the pretrained model.
\item[\textup{(iii)}] The conditions \eqref{eq:nonreach-cond} hold for
$\beta_h=\beta_f=1$, $\beta_v=0$, $\Delta R=C=1$, $\kappa=0.1$,
$q_0=s_0=0.01$, $\eta=20$ (then $A_0=-1$, $D=2$,
$A_0+Ds_0=-0.98$, $Dq_0=0.02$, $A_0+D-\kappa=0.9>\frac{1}{20}\log10^4\approx0.4605$).
\end{enumerate}
\end{proposition}

\begin{proof}
(i) The first two conditions are $f(x_0)\le0$, so
Corollary~\ref{cor:penalty-window}(iii) and
Theorem~\ref{thm:dynamics-two-channel}(c2) apply. For $q\le q_0$ and
$s\le s_0$ we have $A_0+Ds\le A_0+Ds_0\le0$, hence
$J(q,s)\le0$ by \eqref{eq:two-channel-J}. (ii) $J$ is continuous on
$[0,1]^2$ (with $0\log0=0$) and
$J(1,1)=A_0+D-\kappa-\eta^{-1}\log\frac1{q_0s_0}>0$. For fixed $s$,
$q\mapsto J(q,s)$ is strictly concave with derivative
$A_0+Ds-\eta^{-1}(\operatorname{logit}q-\operatorname{logit}q_0)$, which
tends to $+\infty$ as $q\downarrow0$ and to $-\infty$ as $q\uparrow1$, so
its maximum is interior; the same holds in $s$ for fixed $q$. A global
maximiser $(q^\star,s^\star)$ maximises $J$ in each variable separately,
hence is interior and satisfies $\partial_qJ=\partial_sJ=0$, i.e.\ it is an
equilibrium $x^\star$, with $J(x^\star)\ge J(1,1)>0$. By (i),
$x^\star\ne x_\star$ and $x^\star\not\le x_0$. By
Theorem~\ref{thm:dynamics-two-channel}(b1) either $x^\star<x_\star$, which
would give $x^\star\le x_0$, or $x^\star>x_\star$. (iii) is arithmetic.
\end{proof}

In the instance (iii), a numerical solution of $\Phi(w)=w$ gives three
equilibria, at $(q,s)\approx(2.2\cdot10^{-11},\,1.37\cdot10^{-3})$,
$(0.173,\,0.576)$ and $(1-2\cdot10^{-7},\,1-3\cdot10^{-15})$, with
$J\approx-9.4\cdot10^{-4}$, $-0.147$ and $0.4395$ respectively; the
conditions \eqref{eq:cond-L}--\eqref{eq:cond-H} also hold.

\paragraph{Operator~A in the two-channel model.}
Assume $\beta_v>0$. Under operator~A the effective reward is
$R+\eta^{-1}\log\beta$, that is, $\eta^{-1}\log\beta_h$ for an honest
trajectory, $\Delta R+\eta^{-1}\log\beta_v$ for a transparent hack and
$\Delta R+\eta^{-1}\log\beta_f$ for an obfuscated hack. Up to an additive
constant the objective on product policies is again of the form
\eqref{eq:two-channel-J}, with
\[
A_0^{\mathsf A}:=\Delta R-\eta^{-1}\log\frac{\beta_h}{\beta_v},\qquad
D^{\mathsf A}:=\eta^{-1}\ell_v,\qquad \ell_v=\log\frac{\beta_f}{\beta_v}>0,
\]
in place of $(A_0,D)$, so Theorem~\ref{thm:dynamics-two-channel} applies
verbatim. The difference is that $\eta D^{\mathsf A}=\ell_v$ does not
depend on $\eta$ or on any penalty.

\begin{proposition}[Two-channel dynamics under operator~A]
\label{prop:two-channel-A}
Assume $\beta_v>0$ and consider \eqref{eq:two-channel-flow} with
$(A_0^{\mathsf A},D^{\mathsf A})$.
\begin{enumerate}
\item[\textup{(i)}] If $\beta_f/\beta_v<e^{4}$, then for every $\eta>0$ and
$\kappa\ge0$ there is a unique equilibrium, and every solution converges
to it with
$\|x(t)-x^\ast\|_\infty\le e^{-(1-\ell_v/4)t/\eta}\|x(0)-x^\ast\|_\infty$.
\item[\textup{(ii)}] Every equilibrium satisfies
$w^\ast=w_0-\eta\kappa+\ell_v\,q^\ast\le w_0+\ell_v$ and
$u^\ast=u_0+\eta\Delta R-\log(\beta_h/\beta_v)+\ell_v\,s^\ast
\ge u_0+\eta\Delta R-\log(\beta_h/\beta_v)$. Hence the odds of obfuscation
at equilibrium are at most $\beta_f/\beta_v$ times the pretrained odds,
whatever $\eta$, while $q^\ast\to1$ as $\eta\to\infty$.
\item[\textup{(iii)}] The solution started at $x_0$ satisfies
$q(t)\le q_0$ and $s(t)\le s_0$ for all $t$ if and only if
$\eta\Delta R\le\log(\beta_h/\beta_v)-s_0\ell_v$ and
$q_0\ell_v\le\eta\kappa$.
\end{enumerate}
\end{proposition}

\begin{proof}
(i) is Theorem~\ref{thm:dynamics-two-channel}(b4), since
$\eta D^{\mathsf A}=\ell_v<4$ and
$1/\eta-D^{\mathsf A}/4=(1-\ell_v/4)/\eta$. (ii) is the pair of equations
$u^\ast=U(w^\ast)$, $w^\ast=W(u^\ast)$ with
$\eta A_0^{\mathsf A}=\eta\Delta R-\log(\beta_h/\beta_v)$ and
$\eta D^{\mathsf A}=\ell_v$. (iii) is
Corollary~\ref{cor:penalty-window}(iii)(1)$\Leftrightarrow$(2).
\end{proof}

\begin{remark}[Basins under A and under B]\label{rem:A-vs-B}
Under operator~B the coupling $\eta D=\eta(\beta_f-\beta_v)(\Delta R+C)$
grows with the pressure and with the penalty, and
Proposition~\ref{prop:bistability} gives two basins for large $\eta$ as
soon as transparent hacking is deterred ($A_0<0$) and obfuscated hacking
is profitable ($A_0+D>0$): the outcome then depends on the initial
condition, the honest--transparent basin can contain the pretrained model
at arbitrarily large pressure
(Corollary~\ref{cor:penalty-window}), and the equilibrium reached need not
be the maximiser of the objective (Proposition~\ref{prop:nonreach}). Under
operator~A the coupling is the constant $\log(\beta_f/\beta_v)$: unless
the monitor separates the two kinds of hack by a factor of more than
$e^4\approx54.6$ there is a single basin and no history dependence; the
amplification of obfuscation is capped by the factor $\beta_f/\beta_v$,
but hacking is not deterred at large pressure, in agreement with the
static threshold. We have not analysed operator~A for
$\beta_f/\beta_v\ge e^4$ beyond the general statements of
Theorem~\ref{thm:dynamics-two-channel}, but the threshold is not vacuous: for $\beta=(1,10^{-3},0.5)$, $\Delta R=1$, $\eta=6$, $q_0=s_0=0.1$ and $\kappa=0.3$ there are three equilibria, $(0.050,0.024)$, $(0.788,0.711)$ and $(0.866,0.799)$, and three equilibria occur in $73$ of $3000$ random draws with $\beta_v=10^{-3}$.
\end{remark}

\begin{remark}[Bifurcations]\label{rem:saddle-node}
By the implicit function theorem, a hyperbolic equilibrium persists and
depends smoothly on the parameters $(\eta,C,\kappa,\dots)$, and all
equilibria lie in the bounded rectangle of
Theorem~\ref{thm:dynamics-two-channel}(a). Hence the number of equilibria
is locally constant on the open set of parameters at which every
equilibrium is hyperbolic, and can change only at parameters for which
$\Phi$ has a fixed point with $\Phi'=1$. By (b2), such a tangency is a fold
(two equilibria merge, the generic saddle--node scenario) unless it occurs
at the inflection point $w^\dagger$, where the three equilibria merge. We
do not verify the transversality conditions of the saddle--node theorem
along specific paths in the $(\eta,C)$ plane; the existence of the
bistable region itself is Proposition~\ref{prop:bistability}, and
uniqueness for $\eta D<4$ is (b4).
\end{remark}

\begin{remark}[Factorisation is what creates the second basin]
\label{rem:mean-field}
If the policy were tabular over the four joint outcomes (hack or not,
obfuscated or not), Theorem~\ref{thm:dynamics-tabular}(a) would apply with
rewards $0$, $-\kappa$, $A_0$ and $A_0+D-\kappa$: the trajectory would be
the unique Gibbs path and would converge to the unique maximiser, with no
bistability. The two basins of Theorem~\ref{thm:dynamics-two-channel}
arise because the two channels have separate parameters, so that the
propensity to obfuscate is shared between honest and hacking episodes and
the objective restricted to product policies is not concave. Which of the
two idealisations is closer to a given trained model is an empirical
question that this section does not address.
\end{remark}

\begin{remark}[Time scales and noise]\label{rem:noise}
Parts (b4) and (d) give the time to reach an $\varepsilon$-neighbourhood of
a stable node: of order $(1/\eta-D/4)^{-1}\log(1/\varepsilon)$ globally
when $\eta D<4$, and asymptotically $(\mu^\ast)^{-1}\log(1/\varepsilon)$ in
general; at the two nodes of the large-$\eta$ bistable regime
$\Phi'(w^\ast)$ is close to $0$ and $\mu^\ast$ is close to $1/\eta$, the
relaxation rate of the anchor. With stochastic (finite-batch) gradients,
the order certificates of Theorem~\ref{thm:dynamics-two-channel}(c) no
longer hold pathwise, and a transition from the honest--transparent basin
to the hack--obfuscated one becomes a noise-activated event, whose
frequency should depend on the batch size, the step size and the height
of the barrier $\hat J(x_1)-\hat J(x_2)$ in the Fisher geometry. We make no
quantitative claim about this rate.
\end{remark}

\begin{remark}[What is known and what is specific to this section]
\label{rem:known-vs-new}
The mathematical tools and several of the models are standard. The multiplicative form and linear
convergence of entropy-regularised natural policy gradient are due to
\cite{cen2022fast} (see also \cite{agarwal2021theory} for the
unregularised case), and convergence rates of softmax policy gradient to
\cite{mei2020global}; Theorem~\ref{thm:dynamics-tabular}(a) is a direct
consequence of the known update in the bandit case. The order-theoretic
arguments in Theorem~\ref{thm:dynamics-two-channel}(c) are elementary
instances of the theory of cooperative systems
\cite{smith1995monotone,hirsch2006monotone}, proved here directly, and the
convergence arguments use LaSalle's invariance principle
\cite{lasalle1976stability}. The tabular natural-gradient lemma goes back to \cite{kakade2001natural}. The two-channel system is the classical pair of leaky, mutually exciting sigmoid units (continuous Hopfield and Wilson--Cowan models, with the Cohen--Grossberg--Hopfield Lyapunov function~\cite{hopfield1984,wilson1972,cohen1983}); bistability, at most three equilibria and fold bifurcations are classical for it. It is also regularised logit-response dynamics in a $2\times2$ potential game, whose temperature bifurcations are studied in~\cite{hofbauer2002,leonardos2022}, and the non-concavity of the objective restricted to product policies is the classical mean-field phenomenon~\cite{wainwright2008}. What is specific to this section is only the
application of these tools and models to in-loop filters during post-training and the comparison of learning rules in their presence: the time-slice form of operators~A and~B and the resulting
exact dominance times; the distinction between distribution-matching and
advantage-based readings of rejection
(Theorem~\ref{thm:dynamics-tabular}(b)--(c),
Proposition~\ref{prop:discard-kl}), including the identification of the
discard policy gradient with ascent on the accepted-sample mean reward;
and the two-channel analysis (count and ordering of equilibria,
certificates, penalty windows, non-reachability). We have not conducted a
systematic literature search establishing that these specific statements
are absent from prior work.
\end{remark}
